\labsection{8}{Testing whether a cryptocurrency is money}{sec:lab08-crypto}

\begin{labproject}{Apply the three-function test with evidence rather than opinion}
\textbf{Coverage.} Chapter~\ref{ch:cryptocurrencies}.\\
\textbf{Goal.} Measure the properties Chapter~\ref{ch:cryptocurrencies} discusses qualitatively, and reach a defensible verdict on each of money's three functions.

\courselabflow{Load prices}{Compute returns}{Measure drawdown}{Assess functions}



\end{labproject}

\textbf{Complete worked example.} This runnable fixture uses synthetic inputs. First reproduce it, then adapt the data or ratings with documented evidence. All code below belongs to one file.
\begin{lstlisting}[style=utpterminal]
python all_experiments/lab08_crypto_as_money.py
\end{lstlisting}

\noindent\textbf{Setup and output folder.}
\lstinputlisting[style=utppython,firstline=1,lastline=10]{all_experiments/lab08_crypto_as_money.py}

\noindent\textbf{Generate comparable daily teaching series.}
\lstinputlisting[style=utppython,firstline=11,lastline=23]{all_experiments/lab08_crypto_as_money.py}

\noindent\textbf{Measure volatility on the same observation frequency.}
\courselabtrace{Peak-to-trough loss is not daily volatility}{Peak}{Maintain the largest price observed so far}{Drawdown}{Current price divided by running peak, minus one}{Worst loss}{The minimum of that series is maximum drawdown}



\textbf{Use the chapter while coding.} Revisit \hyperref[ch:cryptocurrencies]{the earlier chapter explanation}; match its operation to the variables and state changes below.

\textbf{Follow the supplied implementation.} The frequency factor belongs to this daily fixture. Inspect the dates and missing values before reusing it for another series. Recovery is the gain from the trough back to the prior peak, not the negative of the drawdown percentage.
\courselabtrace{Trace the synthetic daily risk calculation}{1. Input / state}{px has 366 daily prices; log(px / px.shift(1)).dropna() yields 365 returns.}{2. Operation}{annualised uses sample standard deviation times sqrt(365); drawdowns divides prices by the running maximum.}{3. Check}{worst is each column's minimum drawdown; recovery = 1 / (1 + worst) - 1.}


\lstinputlisting[style=utppython,firstline=24,lastline=30]{all_experiments/lab08_crypto_as_money.py}

\noindent\textbf{Measure path-dependent drawdown and recovery.}
\lstinputlisting[style=utppython,firstline=31,lastline=39]{all_experiments/lab08_crypto_as_money.py}

\noindent\textbf{Compare explicit payment assumptions.}
\lstinputlisting[style=utppython,firstline=40,lastline=45]{all_experiments/lab08_crypto_as_money.py}
\begin{labvisual}{Synthetic series: largest observed peak-to-trough loss}
\begin{labgraphic}\begin{tikzpicture}[font=\small]
\node[anchor=east,font=\scriptsize] at (-.2,0.0) {high volatility};
\fill[UTPPaleBlue] (0,-0.16) rectangle (7,0.16);
\fill[UTPTeal] (0,-0.16) rectangle (3.6685,0.16);
\node[anchor=west,font=\scriptsize] at (7.15,0.0) {52.4\% loss};
\node[anchor=east,font=\scriptsize] at (-.2,-0.55) {medium volatility};
\fill[UTPPaleBlue] (0,-0.7100000000000001) rectangle (7,-0.39);
\fill[UTPTeal] (0,-0.7100000000000001) rectangle (2.5959,-0.39);
\node[anchor=west,font=\scriptsize] at (7.15,-0.55) {37.1\% loss};
\node[anchor=east,font=\scriptsize] at (-.2,-1.1) {low volatility};
\fill[UTPPaleBlue] (0,-1.26) rectangle (7,-0.9400000000000001);
\fill[UTPTeal] (0,-1.26) rectangle (0.5121,-0.9400000000000001);
\node[anchor=west,font=\scriptsize] at (7.15,-1.1) {7.3\% loss};
\end{tikzpicture}\end{labgraphic}\end{labvisual}

\textbf{Check the run.} The following output was captured from the supplied script; it describes the synthetic fixture only.
\lstinputlisting[style=utpterminal]{labs/outputs/lab08.txt}
% Reused from chapters/08-cryptocurrencies.tex
\begin{labvisual}{Recall the chapter model}
\begin{labgraphic}
\begin{tikzpicture}[
  p/.style={draw=UTPBlue,fill=UTPPaleBlue,rounded corners=2mm,text width=2.5cm,minimum height=1.0cm,align=center,font=\scriptsize\bfseries\color{UTPNavy}},
  net/.style={draw=UTPTeal,fill=UTPPaleTeal,rounded corners=2mm,text width=3.1cm,minimum height=1.3cm,align=center,font=\scriptsize\bfseries\color{UTPNavy}},
  arr/.style={-{Stealth[length=2mm]},thick,draw=UTPTeal},
  bad/.style={-{Stealth[length=2mm]},thick,dashed,draw=UTPRed}
]
\node[p] (a) at (-4.5,0) {Alice};
\node[net] (n) at (0,0) {Shared network\\validates spend history};
\node[p] (b) at (4.5,1.2) {Bob}; \node[p] (c) at (4.5,-1.2) {Charlie};
\draw[arr](a)--node[above,font=\scriptsize]{signed spend}(n); \draw[arr](n)--(b);
\draw[bad](a) to[bend right=18] node[below,font=\scriptsize]{conflicting second spend?}(n);
\draw[bad](n)--(c);
\end{tikzpicture}
\end{labgraphic}\end{labvisual}


\begin{verification}
\textbf{Explain and extend your result.} Identify the largest synthetic drawdown and calculate the recovery needed. Compare each fee with a small purchase amount and state which assumption most affects payment practicality.

\begin{itemize}\item The script runs from a fresh process and creates the named CSV or report files.\item Every rate uses a stated population and denominator; every model assumption is identified.\item Synthetic outputs are labelled and are not presented as observations of a real institution.\item Change one input, rerun, and explain the direction of the result.\end{itemize}\end{verification}
\begin{labdeliverable}Submit the complete script, generated evidence, and a one-page interpretation linking the results to the chapter model. State the assumptions and one limitation.\end{labdeliverable}
